\chapter{Results}
\label{ch:results}

\section{Interference}
\label{sec:interference}

Now that the B210 has been configured for radar testing it is ready to be used.
However, in initial testing a target could not be detected.
The presence of the feed through leakage, signal directly from the transmitter to the receiver,
was evident, but no other targets could be seen.
The gain was adjusted, wired connections checked, signal processing improved,
but there was still no visible targets.
Tests were preformed to see how the system functioned, testing individual channels,
terminating connections, changing gain values.
The SDR functioned as expected, however, when transmit and receive
were active at the same time, there was immense leakage
even when the outputs were terminated.
The interferance is internal to the SDR and is powerful enough to drown at potential targets.
The high interferance only presents issues with close targets as they are drowned out.
Since testing the radar involves close-range targets it is necessary to identify and resolve the issue.
In detecting avalanches this presents less of an issue since
the radar does not need to identify close-range targets, only targets hundreds or thousands of meters away.
To test the interferance, the transmit and receive on one channel were terminated.
The gain was set to half to measure the intensity of the internal leakage.
The results are shown in Figure \ref{fig:half_power_leakage}.
At low power the internal interferance isn't significant but as transmit power increases
the interferance quickly dominates as shown in Figure \ref{fig:three_fourths_power_leakage}.

\begin{figure}[h]
    \includegraphics[width=1\textwidth, center]{internal_interferance_1_2.png}
    \caption{Return power from internal leakage. Transmit and receive ends are terminated.
            The power at 0m shows the internal leakage.
            Gain in dB is set to about half of the total available.}
    \label{fig:half_power_leakage}

    \includegraphics[width=1\textwidth, center]{internal_interferance_3_4.png}
    \caption{Return power from internal leakage. Transmit and receive ends are terminated.
            The power at 0m shows the internal leakage.
            Gain in dB is set to about three fourths of the total available.
            The noise from the leakage is so severe that close targets can't be detected.}
    \label{fig:three_fourths_power_leakage}
\end{figure}

The internal leakage makes it difficult to detect targets.
To combat this, it was tested using one channel to receive
and the other to transmit.
By seperating the channels the interferance dropped as is evident when comparing
Figures \ref{fig:three_fourths_power_leakage} and \ref{fig:two_channel_power_leakage},
both datasets were taken at the same gains.
The reduced interferance is enough to successfully capture and view radar targets.

\begin{figure}[h]
    \includegraphics[width=1\textwidth, center]{internal_interferance_two_3_4.png}
    \caption{Return power from internal leakage. Transmit and receive ends are terminated.
            The power at 0m shows the internal leakage.
            Gain in dB is set to about three fourths of the total available.
            The frequency bins are wider because of a lower sample rate (30 MSPS).
            In this case channel 0 is used for transmit and channel 1 is used for receive
            to attempt to reduce internal leakage.
            The Figure shows a dramatic decrease in interferance
            due to the sepeartion of the two channels.}
    \label{fig:two_channel_power_leakage}
\end{figure}

% Explain the results I found and what they mean
\section{Radar Testing}
With the radar working with limited interferance the system as a whole could be verified.
Unfortunately, two channel operation can not be performed without reintroducing
the high levels of interferance from transmitting and receinving on the same channel.
However, the higher sample rate of 56 MSPS can be used as only one receive channel is used.
Tests were performed in an open courtyard with a few trees and lamposts as obstrutions.
The SDR was set to collect samples as a corner reflector was moved away from the radar.
The results are shown in Figure \ref{fig:corner_moving_away}.
The radar successfully detected the corner reflector as it moved away over time.


\begin{figure}[h]
    \includegraphics[width=1\textwidth, center]{corner_moving_away.png}
    \caption{Return power plotted against time and reflection distance.
                The corner reflector is seen as a bright spot
                increasing in distance over time.
                The reflector is easily seen travelling out to around 80 meters.}
    \label{fig:corner_moving_away}
\end{figure}
